%%%%%%%%%%  scan idx 438  |  folio 425  %%%%%%%%%%
le sous-schéma fermé de \(X\) défini par \(J\), \(X'\) le schéma obtenu en faisant éclater \(Y\), \(P\) le fibré projectif associé à \(E\), \(p\) le morphisme structural \(P\to X\). Sur \(P\), on considère la suite exacte canonique :
\[
0\to F\to p^*(E)\to\cO_P(1)\to0,
\]
et le morphisme composé :
\[
 v:F\longrightarrow p^*(E)\xrightarrow{p^*(u)}\cO_P .
\]

Alors :
\begin{enumerate}[label=\roman*)]
\item La formation de \(X'\) commute à tout changement de base \(f:T\to X\) tel que \(f^*(u)\) soit régulier.
\item L'immersion fermée \(j:X'\to P\) définie par l'homomorphisme surjectif \(\mathbf{S}(E)\longrightarrow\bigoplus_{n\ge0}J^n\) défini par \(u\) en degré \(1\) est une immersion régulière, son idéal étant l'image du morphisme régulier \(v\) ; en particulier si \(\operatorname{codim}(i)\) est constante, \(\operatorname{codim}(j)=\operatorname{codim}(i)-1\).
\item Le \(\mathbf{S}(E)\)-Module \(\bigoplus_{n\ge0}J^n\) est parfait en tant que \(\mathbf{S}(E)\)-Module gradué, i.e. il existe localement une résolution finie, graduée, de \(\bigoplus_{n\ge0}J^n\) par des \(\mathbf{S}(E)\)-Modules gradués localement libres de type fini (VI 1.2).
\end{enumerate}

i) Comme \(X'=\Proj(\bigoplus_{n\ge0}J^n)\), le résultat est une conséquence immédiate des relations (1.2.3) et de EGA II.

ii) Les propriétés énoncées étant locales sur \(X\), on peut supposer \(X\) affine, et de plus \(E\) libre sur \(X\). On pose \(X=\Spec(A)\).

On fait d'abord la démonstration dans le cas standard \(A=\mathbb Z[T_1,\ldots,T_d]\), \(E\) étant muni d'une base \(e_1,\ldots,e_d\), avec \(u(e_i)=T_i\). Alors \(P=\Proj\bigl(\mathbb Z[T_1,\ldots,T_d,X_1,\ldots,X_d]\bigr)\), la graduation étant prise par rapport aux \(X_i\). Pour montrer que \(v\) est un morphisme régulier dont l'image est l'Idéal de \(j\), il suffit de le faire localement sur \(P\), donc sur les ouverts
